\scanpage{19}
\mspage{9}
La notion d'isotopie raisonnable pour des
tels morphismes (dans le contexte topologique)
se trouve être équivalente à la congruence
sous le groupe $\operatorname{Aut}^0(X)\times\operatorname{Aut}^0(Y)$
(où $\operatorname{Aut}^{\deleted{\unsure{0}}}(X)$ est interprété comme $\operatorname{Aut}(\widehat X,S)$,
et $\operatorname{Aut}(Y)$ comme $\operatorname{Aut}(\widehat{\unsure{X|Y}},T)$). L'ennui,
c'est qu'il n'y a plus de factorisation
\uline{canonique} d'une classe d'isotopie des
morphismes \inserted{$X\longrightarrow Y$} (dont les éléments \inserted{\unsure{variables}} $f$
transforment $S_1=S\setminus S_0$ où $S_0=\widehat X|T$) en
des parties variables de $Y=\widehat Y\setminus T$) en
un isomorphisme \deleted{suivi} isotopique \deleted{suivi} \inserted{ou plutôt}
\inserted{une classe d'isotopie d'inclusion type ``bouchage de trous'',}
\inserted{suivi} d'un morphisme $Y'\longrightarrow Y$, où $Y'$ est un
morphisme fini (possédant maintenant aussi
\inserted{en plus} des pts de ramification, au dessus des
points de $Y$ qui \uline{ne sont pas} déterminés
par la classe seule). Pour cette raison,
je ne vois plus comment interpréter une
telle classe, ou son image dans un
\deleted{``\unsure{complété} profini''} ``complété profini'' convenable,
en termes de géom. alg. ``abstraite'' sur $\mathbb C$,
\margin{18}
